\subsection{Acertijo lógico}

\subsubsection{Descripción del problema}

El problema consiste en determinar el apellido, la edad y el género musical preferido de tres personas: Juan, Darío y Óscar. Los apellidos posibles son López, Gonzales y García; las edades son 24, 25 y 26 años; y los géneros musicales son Clásica, Pop y Jazz.

El objetivo es encontrar una asignación que cumpla todas las pistas del acertijo, garantizando que cada persona tenga una edad, un apellido y un género musical diferente.

\subsubsection{Modelo}

\paragraph{Parámetros}

Se definen los siguientes conjuntos enumerados:

\begin{itemize}
    \item Personas:
    \[
    Nombres = \{Juan, Dario, Oscar\}
    \]
    \item Apellidos:
    \[
    Apellidos = \{Lopez, Gonzales, Garcia\}
    \]
    \item Géneros musicales:
    \[
    TipoMusica = \{Clasica, Pop, Jazz\}
    \]
\end{itemize}

El dominio de edades se define como:

\[
Edades = \{24,25,26\}
\]

\paragraph{Variables y dominios}

Se definen tres arreglos de variables de decisión:

\begin{itemize}
    \item Edad:
    \[
    edad[i] \in \{24,25,26\}
    \quad \forall i \in Nombres
    \]

    \item Apellido:
    \[
    apellido[i] \in Apellidos
    \quad \forall i \in Nombres
    \]

    \item Música:
    \[
    musica[i] \in TipoMusica
    \quad \forall i \in Nombres
    \]
\end{itemize}

En total, el modelo contiene nueve variables de decisión, tres por cada característica.

\paragraph{Restricciones}

Las restricciones del problema se dividen en condiciones de unicidad y pistas específicas.

\textbf{1. Unicidad de valores}

Cada persona debe tener una edad, un apellido y un género musical diferentes. El enunciado no dice explícitamente que las edades sean distintas, solo que están entre 24 y 26 años; se asume porque, sin esa condición, el acertijo tiene cinco soluciones en lugar de una. Se expresa mediante:

\[
\forall i,j \in Nombres,\ i \neq j:
\]

\[
edad[i] \neq edad[j]
\]

\[
apellido[i] \neq apellido[j]
\]

\[
musica[i] \neq musica[j]
\]

Estas condiciones se implementan mediante desigualdades explícitas o utilizando la restricción global \texttt{alldifferent}, según la implementación seleccionada.

\textbf{2. Apellido Gonzales}

La persona con apellido Gonzales debe ser menor que Juan y preferir música clásica:

\[
\forall j \in Nombres:\;
apellido[j] = Gonzales
\Rightarrow
\left(
edad[Juan] > edad[j]
\land
musica[j] = Clasica
\right)
\]

\textbf{3. Música Pop}

Quien prefiera música Pop no puede tener el apellido García ni tener 24 años:

\[
\forall i \in Nombres:\;
musica[i] = Pop
\Rightarrow
\left(
apellido[i] \neq Garcia
\land
edad[i] \neq 24
\right)
\]

\textbf{4. Condiciones de Óscar}

Óscar no tiene el apellido López y tiene 25 años:

\[
apellido[Oscar] \neq Lopez
\]

\[
edad[Oscar] = 25
\]

\textbf{5. Condición de Darío}

Darío no prefiere música Jazz:

\[
musica[Dario] \neq Jazz
\]

\paragraph{Restricciones redundantes}

La tercera implementación incorpora dos restricciones redundantes que se deducen de las pistas originales:

\[
edad[Juan] \neq 24
\]

\[
musica[Juan] \neq Clasica
\]

La primera se deduce porque Óscar tiene 25 años y todas las edades deben ser diferentes. Además, la persona Gonzales debe ser menor que Juan. Por lo tanto, Juan no puede tener 24 años y necesariamente debe tener 26.

La segunda se obtiene porque Juan no puede tener el apellido Gonzales: esto implicaría que su edad fuera mayor que sí misma. Como la persona Gonzales prefiere música clásica y los géneros deben ser diferentes, Juan tampoco puede preferirla.

Estas restricciones no modifican las soluciones posibles, sino que ayudan al solver a deducir información anticipadamente.

\paragraph{Ruptura de simetrías}

No se implementan restricciones adicionales de ruptura de simetrías. Las personas tienen identidades específicas y algunas pistas hacen referencia directamente a Juan, Darío y Óscar, lo que permite diferenciarlas dentro del modelo.

\subsubsection{Detalles de implementación}

El modelo se desarrolló en MiniZinc utilizando conjuntos enumerados y arreglos de variables indexados por el conjunto \texttt{Nombres}.

Se definieron dos parámetros para controlar las pruebas:

\begin{itemize}
    \item \texttt{impl}: selecciona la implementación de las restricciones de unicidad.
    \item \texttt{strategy}: selecciona la estrategia de búsqueda.
\end{itemize}

Se implementaron tres versiones:

\begin{itemize}
    \item \textbf{Implementación 1:} utiliza desigualdades explícitas entre cada par de personas.
    \item \textbf{Implementación 2:} utiliza la restricción global \texttt{alldifferent} para edades, apellidos y música.
    \item \textbf{Implementación 3:} utiliza \texttt{alldifferent} y agrega las dos restricciones redundantes.
\end{itemize}

También se construye el arreglo \texttt{vars\_busqueda}, que contiene las nueve variables en el siguiente orden: primero las edades, luego los apellidos y finalmente los géneros musicales.

El modelo utiliza \texttt{solve satisfy}, debido a que el objetivo es encontrar una asignación que cumpla todas las restricciones, sin optimizar una función objetivo.

\subsubsection{Estrategias de búsqueda}

Se implementaron tres estrategias de búsqueda. Todas utilizan \texttt{indomain\_min}, que intenta asignar primero el menor valor disponible del dominio.

\begin{itemize}
    \item \textbf{Estrategia 1: input\_order.} Selecciona las variables siguiendo el orden en que aparecen en \texttt{vars\_busqueda}.

    \item \textbf{Estrategia 2: first\_fail.} Selecciona primero la variable que tenga el dominio actual más pequeño, buscando detectar contradicciones anticipadamente.

    \item \textbf{Estrategia 3: smallest.} Selecciona la variable cuyo menor valor disponible sea el más pequeño.
\end{itemize}

La Tabla~\ref{tab:acertijo-estrategias} resume las estrategias.

\begin{table}[htbp]
\caption{Estrategias de búsqueda implementadas}
\label{tab:acertijo-estrategias}
\begin{tabular}{@{}ccc@{}}
\toprule
\textbf{Estrategia} & \textbf{Selección de variable} & \textbf{Asignación de valor} \\
\midrule
1 & input\_order & indomain\_min \\
2 & first\_fail & indomain\_min \\
3 & smallest & indomain\_min \\
\bottomrule
\end{tabular}
\end{table}

\subsubsection{Pruebas}

Las pruebas se realizaron utilizando la instancia \texttt{acertijo\_01} y el solver Gecode, combinando las tres implementaciones con las tres estrategias de búsqueda.

La configuración utilizada fue:

\begin{itemize}
    \item \textbf{Solver:} Gecode 6.3.0.
    \item \textbf{MiniZinc:} versión 2.9.3.
    \item \textbf{Tiempo máximo:} 60 segundos.
    \item \textbf{Instancias:} acertijo\_01.
    \item \textbf{Implementaciones y estrategias:} tres de cada una.
\end{itemize}

Se utilizó la opción \texttt{-s} de MiniZinc para obtener estadísticas de resolución, incluyendo tiempo, nodos, fallos y profundidad máxima.

La Tabla~\ref{tab:acertijo-resultados} muestra los resultados.

\begin{table}[htbp]
\small
\caption{Resultados de las pruebas con Gecode}
\label{tab:acertijo-resultados}
\begin{tabular}{@{}ccccccc@{}}
\toprule
\textbf{Impl.} & \textbf{Estrategia} & \textbf{Estado} & \textbf{Tiempo (s)} & \textbf{Nodos} & \textbf{Fallos} & \textbf{Prof.} \\
\midrule
1 & 1 & sat & 0,0001276 & 3 & 1 & 1 \\
1 & 2 & sat & 0,0001284 & 3 & 1 & 1 \\
1 & 3 & sat & 0,0001315 & 5 & 1 & 3 \\
2 & 1 & sat & 0,0001322 & 3 & 1 & 1 \\
2 & 2 & sat & 0,0001418 & 3 & 1 & 1 \\
2 & 3 & sat & 0,0001370 & 5 & 1 & 3 \\
3 & 1 & sat & 0,0000979 & 1 & 0 & 0 \\
3 & 2 & sat & 0,0000953 & 1 & 0 & 0 \\
3 & 3 & sat & 0,0001555 & 1 & 0 & 0 \\
\bottomrule
\end{tabular}
\end{table}

Todas las ejecuciones finalizaron con estado \texttt{sat}, indicando que se encontró una solución válida. No se utilizó la opción \texttt{-a}, por lo que no se enumeraron todas las soluciones posibles.

\subsubsection{Análisis de los árboles de búsqueda}

En las implementaciones 1 y 2 se obtuvieron los mismos resultados en cuanto a nodos, fallos y profundidad (Tabla~\ref{tab:acertijo-arbol}):

\begin{itemize}
    \item Estrategias 1 y 2: 3 nodos, 1 fallo y profundidad 1.
    \item Estrategia 3: 5 nodos, 1 fallo y profundidad 3.
\end{itemize}

Por otra parte, la implementación 3 obtuvo 1 nodo, 0 fallos y profundidad 0 en las tres estrategias.

Esto muestra que las restricciones redundantes facilitaron la propagación, permitiendo encontrar una solución sin necesidad de realizar ramificaciones registradas en las estadísticas.

\begin{table}[htbp]
\caption{Estadísticas de los árboles de búsqueda}
\label{tab:acertijo-arbol}
\begin{tabular}{@{}ccccc@{}}
\toprule
\textbf{Implementación} & \textbf{Estrategia} & \textbf{Nodos} & \textbf{Fallos} & \textbf{Profundidad} \\
\midrule
1 & 1 & 3 & 1 & 1 \\
1 & 2 & 3 & 1 & 1 \\
1 & 3 & 5 & 1 & 3 \\
2 & 1 & 3 & 1 & 1 \\
2 & 2 & 3 & 1 & 1 \\
2 & 3 & 5 & 1 & 3 \\
3 & 1 & 1 & 0 & 0 \\
3 & 2 & 1 & 0 & 0 \\
3 & 3 & 1 & 0 & 0 \\
\bottomrule
\end{tabular}
\end{table}

Estos resultados corresponden a búsquedas detenidas al encontrar una solución, por lo que no representan necesariamente el árbol completo.

\subsubsection{Análisis de resultados}

Las implementaciones 1 y 2 presentaron estadísticas de búsqueda idénticas. Esto indica que, para esta instancia, utilizar \texttt{alldifferent} en lugar de desigualdades explícitas no produjo una mejora observable en la cantidad de nodos ni en los fallos.

En cambio, la implementación 3 obtuvo mejores resultados en estos indicadores. Al agregar las restricciones redundantes, el solver pudo deducir anticipadamente información de las pistas originales, reduciendo los nodos de 3 o 5 a solamente 1 y eliminando los fallos registrados.

En cuanto al tiempo de búsqueda, el menor valor fue de 0,0000953 segundos, correspondiente a la implementación 3 con la estrategia 2. El mayor fue de 0,0001555 segundos, correspondiente a la implementación 3 con la estrategia 3.

Los tiempos son extremadamente pequeños, por lo que sus diferencias deben interpretarse con precaución. Además, el tiempo de compilación estuvo entre 0,176 y 0,189 segundos, considerablemente por encima del tiempo de búsqueda.

Finalmente, dado que solo se evaluó una instancia, no es posible generalizar estos resultados a todos los problemas de acertijos lógicos.

\subsubsection{Conclusiones del ejercicio}

El ejercicio permitió modelar un acertijo lógico mediante variables, dominios y restricciones, utilizando MiniZinc y el solver Gecode.

La comparación mostró que las implementaciones con desigualdades explícitas y \texttt{alldifferent} tuvieron un comportamiento equivalente en las pruebas realizadas. Sin embargo, agregar restricciones redundantes permitió reducir considerablemente los nodos explorados y eliminar los fallos registrados.

También se observó que las estrategias de búsqueda pueden influir en el proceso de resolución, aunque su efecto fue menos relevante en la tercera implementación.

En conclusión, el ejercicio demuestra que incorporar consecuencias lógicas de las restricciones originales puede mejorar la propagación y reducir el esfuerzo de búsqueda. No obstante, los resultados deben considerarse dentro del contexto de la única instancia evaluada.
